\documentclass[mathserif, fleqn]{beamer}
% basic text processing/formatting
\usepackage{xltxtra}            % XeTeX packages
\usepackage[english]{babel}

% maths stuff
\usepackage{amsmath}
\usepackage{amsfonts}
\usepackage{amssymb}
\usepackage{mathtools}
\usepackage{slashed}            % (Feynman-) slashed symbols
\usepackage{commath}            % for nicer differentials
\usepackage{dsfont}             % double stroked characters

\usepackage{tikz}               % drawing all manners of things
\usepackage{forest}             % draw trees using TikZ
\usepackage{datetime}           % for custom date and time format
\usepackage{array}              % for matrices in math environment
\usepackage{float}              % place graphics with "H"
\usepackage{braket}             % Dirac notation
\usepackage{placeins}           % FloatBarrier
\usepackage{lmodern}            % to scale font to whatever I want
\usepackage{verbatim}           % print input as it is
\usepackage{xfrac}              % sideways fractions
\usepackage{ccicons}            % creative commons icons

% For textblock, remove showboxes for release.
% Gets loaded by aiphi, but this way we can specify showboxes.
\usepackage[absolute,overlay,showboxes]{textpos}

% load customisations
\usepackage{aiphi}

%% devel only
% overlay a grid over the slides
% \usepackage[texcoord,grid,gridcolor=black!10,subgridcolor=black!5,gridunit=mm]{eso-pic}
\usepackage{lipsum}


% ---------------------------------------------------------------------------------------
% general info
% ---------------------------------------------------------------------------------------

% commands so we can refer to these later
\newcommand{\theauthor}{Jan-Lukas Wynen}
\newcommand{\theinstitute}{Institute for applied \LaTeX}
\newcommand{\thetitle}{Fancy shmancy talk}
\newcommand{\thesubtitle}{Does this subtitle appear anywhere?}
\newcommand{\thedate}{\isodate\today}
\newcommand{\thelogo}{\includegraphics[width=10mm]{logo.pdf}}
\newcommand{\thelicense}{\href{http://creativecommons.org/licenses/by/4.0/}{\ccby}}

% show just the author on the titlepage but include the license in the footline
\author[\thelicense\quad\theauthor]{\theauthor}
\title{\thetitle}
\date{\thedate}
\institute{\theinstitute}
\logo{\thelogo}


% ---------------------------------------------------------------------------------------
% setup beamer
% ---------------------------------------------------------------------------------------
\usetheme[titleinfoot, toc, rainbow]{aiphi}


% ---------------------------------------------------------------------------------------
% TikZ
% ---------------------------------------------------------------------------------------
\usetikzlibrary{tikzmark, calc, arrows, spy}


% ---------------------------------------------------------------------------------------
% abbreviations
% ---------------------------------------------------------------------------------------
\newcommand{\unit}[1]{\,\text{#1}}
\newcommand{\ev}{\,\text{eV}}
\newcommand{\kev}{\,\text{keV}}
\newcommand{\mev}{\,\text{MeV}}
\newcommand{\gev}{\,\text{GeV}}

\DeclareMathOperator{\Tr}{Tr}

%%% Local Variables:
%%% coding: utf-8
%%% mode: latex
%%% TeX-engine: xetex
%%% TeX-master: "main"
%%% End:


\begin{document}

\begin{frame}
  \titlepage
\end{frame}

\section{A section}
\begin{frame}\frametitle{Overview}
  % different bullets just for this slide
  \setbeamertemplate{itemize item}{}
  \setbeamertemplate{itemize subitem}{\adfrightarrowhead}
  \setbeamertemplate{itemize subsubitem}{--}

  \begin{itemize}\setlength{\itemsep}{1.5ex}
  \item \textbf{Motivation}
    \begin{itemize}\setlength{\itemsep}{1ex}
    \item Latex is `fun'
    \item Really nice slides if you put really much effort into it
    \end{itemize}
  \item \textbf{Method}
    \begin{itemize}\setlength{\itemsep}{1ex}
    \item Editor + PDFLatex
    \item Ah, and PDF viewer
      \begin{itemize}
      \item Many choices
      \item Pretty much all good
      \end{itemize}
    \end{itemize}
  \item \textbf{Results}
    \begin{itemize}\setlength{\itemsep}{1ex}
    \item See for yourself
    \item Looks nice, doesn't it?
    \end{itemize}
  \end{itemize}
\end{frame}

\begin{frame}\frametitle{Conclusion}
  \begin{itemize}\setlength{\itemsep}{2ex}
  \item Woah, that was fast!
    \begin{itemize}
    \item Not much to show here anyway
    \end{itemize}
  \item But look, nice serifs in equations (optional... $uv$):
    \begin{align*}
      \mathcal{C} = \frac1{Z} \braket{0|\int\,\mathcal{D}\psi\,e^{\i S}|0}
    \end{align*}
  \end{itemize}
\end{frame}

\section{Second}
\begin{frame}
  \begin{block}{A normal block}
    Some normal, boring text, but COLOURS!
  \end{block}
  \begin{exampleblock}{An example}
    This is how you do it.
  \end{exampleblock}
  \begin{alertblock}{Careful now}
    You should really watch out more!
  \end{alertblock}
\end{frame}

\begin{frame}[fragile]
  \frametitle{Code!}
  \begin{lstlisting}[caption={A test listing}]
// Oh yeah, I did this!
template <typename It,
          template <typename...> class F,
          typename... TParam>
struct foreach {
    /* magick! */
    template <typename... Args>
    static void f(Args && ...args) {
        const char c = 'c';
        _tmp_internal::foreach_impl<typename It::Head,
                                    typename It::Tail,
                                    F,
                                    TParam...>
                                    ::f("A string", std::forward<Args>(args)...);
                                    std::cout << 'a' << std::endl;
                                    std::vector<int> v(2, 3);
    }
};
  \end{lstlisting}
\end{frame}

\begin{frame}[fragile]
  \lstset{language=python}
  \begin{lstlisting}[caption={Some Python as well}]
# a hat in a cat
class cat:
    def __init__(self, shape):
        self._shape = shape

    def show(self):
        print("A {} shaped hat in a now {} shaped cat".format(self._shape))
  \end{lstlisting}
\end{frame}

\begin{frame}
  Before block
  \begin{textblock*}{200pt} (100pt,1pt)
    \begin{itembox}{first block}
      \begin{itemize}
      \item And-a one
      \item And-a two
      \item And-a
      \end{itemize}
    \end{itembox}
  \end{textblock*}

  \only<2->{
    \begin{textblock*}{200pt} (150pt,50pt)
      \begin{itembox}{second block}
        \begin{itemize}
        \item Foo
        \item Bar
        \item Baz
        \end{itemize}
      \end{itembox}
    \end{textblock*}
  }

  \only<3>{
    \begin{textblock*}{200pt} (200pt,100pt)
      \begin{itembox}{third block}
        \begin{itemize}
        \item En
        \item To
        \item Tre
        \item Fire
        \end{itemize}
      \end{itembox}
    \end{textblock*}
  }
  After block
\end{frame}

\begin{frame}
  \begin{fancybox}{Testing...}
    \lipsum[2]
  \end{fancybox}

  \begin{tagbox}{Title}{Tag}%[aiphired]
    Tagged stuff
  \end{tagbox}
\end{frame}

\begin{frame}[plain,c]
  % suppress strange item marker in center
  \renewcommand*\item{\olditem[]}
  \begin{center}
    \textcolor{aiphired}{\LARGE Thank you for your attention!}
  \end{center}
\end{frame}

\appendix
\backupbegin

\begin{frame}
  \frametitle{Appendix}
  Some extra stuff
\end{frame}

\begin{frame}
  Even more extra stuff
\end{frame}

\backupend

\end{document}